% Source: https://artofproblemsolving.com/community/c3h11311_mathcounts_problems_3
% Stable identity: aops-topic-11311-q02
Symmetry about $P$ means that $P$ is the midpoint of $AR$:
\[
\dfrac{9+a}{2}=5,\qquad\dfrac{14+b}{2}=3\qquad\Longrightarrow\qquad a=1,\ b=-8 .
\]

\begin{asy}
usepackage("amsmath");
size(4.5cm);
defaultpen(fontsize(10pt));
pair A=(9,14),P=(5,3),R=(1,-8);
draw((-1,0)--(11,0),gray,Arrow(TeXHead));draw((0,-10)--(0,16),gray,Arrow(TeXHead));
label("$x$",(11,0),E,gray);label("$y$",(0,16),N,gray);
draw(A--R,blue+linewidth(0.9));
dot(A);dot(P);dot(R,red);
label("$A(9,14)$",A,E);label("$P(5,3)$",P,E);label("$R(1,-8)$",R,E,red);
pair u=unit(rotate(90)*(A-R));
for(pair m: new pair[]{(A+P)/2,(P+R)/2}){draw(m-0.35u--m+0.35u,blue);}
label("\mbox{}",(min(currentpicture,user=true)+max(currentpicture,user=true))/2,black); // keeps the text after the figure black when LaTeX typesets the labels inline
\end{asy}

So $R=\boxed{(1,-8)}$. (Equivalently, $R=2P-A$: from $A$ to $P$ is $(-4,-11)$, and the same step again leads from $P$ to $R$.)
